\chapter{1977年山东A卷}

\section{解答题}

\begin{problem}
\(\sqrt{\left( 3-x \right)^2}=?\)
\end{problem}

\begin{solution}
根号表示算术平方根，因此对任意实数 \(u\)，有 \(\sqrt{u^2}=\lvert u\rvert\). 于是
\(\sqrt{\left( 3-x \right)^2}=\lvert 3-x\rvert\). 当 \(x\le 3\) 时，\(3-x\ge 0\)，所以 \(\lvert 3-x\rvert=3-x\)；当 \(x>3\) 时，\(3-x<0\)，所以 \(\lvert 3-x\rvert=-(3-x)=x-3\). 因而
\[
\sqrt{\left( 3-x \right)^2}=
\begin{cases}
3-x, & x\le 3,\\
x-3, & x>3.
\end{cases}
\]
\end{solution}

\begin{problem}
\(\lg 2+\lg 5+\lg \sqrt{10}+\lg 0.01+\log_8\log_3 3=?\)
\end{problem}

\begin{solution}
因为 \(\lg 2+\lg 5=\lg 10=1\)，\(\lg\sqrt{10}=\frac12\lg 10=\frac12\)，\(\lg 0.01=\lg 10^{-2}=-2\lg 10=-2\)，并且 \(\log_3 3=1\)、\(\log_8 1=0\)，所以原式
\[
\lg 2+\lg 5+\lg\sqrt{10}+\lg 0.01+\log_8\log_3 3=1+\frac12-2+0=-\frac12
\]
\end{solution}

\begin{problem}
写出函数定义域：

\begin{enumerate}
\item \(y=\log_a x\)（\(a>0\)）；
\item \(y=\frac{\sqrt{x-1}}{x+3}\).
\end{enumerate}
\end{problem}

\begin{solution}
\begin{enumerate}
\item 对数的真数必须为正，所以 \(x>0\). 题面给出底数 \(a>0\)；按对数的定义还须有 \(a\ne 1\)，但这不改变 \(x\) 的取值范围. 因此定义域为 \(\{x\mid x>0\}\).
\item 根式有意义要求 \(x-1\ge 0\)，即 \(x\ge 1\)；分母还要求 \(x+3\ne 0\)，即 \(x\ne -3\). 两个条件合并后，\(x\ge 1\) 已自动保证 \(x\ne -3\)，所以定义域为 \(\{x\mid x\ge 1\}\).
\end{enumerate}
\end{solution}

\begin{problem}
写出过点 \((2,3)\)、斜率为 \(-\frac{1}{2}\) 的直线方程.
\end{problem}

\begin{solution}
设所求直线为 \(y=kx+b\). 由斜率条件得 \(k=-\frac12\)，所以直线可写成 \(y=-\frac12x+b\). 将点 \((2,3)\) 代入，得 \(3=-\frac12\cdot2+b\)，从而 \(b=4\). 因此所求直线方程为 \(y=-\frac12x+4\)，等价地也可写成点斜式 \(y-3=-\frac12\left( x-2 \right)\).
\end{solution}

\begin{problem}
求椭圆 \(\frac{x^2}{25}+\frac{y^2}{16}=1\) 的焦点坐标.
\end{problem}

\begin{solution}
将方程写成标准形式
\(\frac{x^2}{5^2}+\frac{y^2}{4^2}=1\). 因为 \(5^2>4^2\)，长轴在 \(x\) 轴上，长半轴 \(a=5\)，短半轴 \(b=4\). 椭圆的焦半距满足 \(c^2=a^2-b^2\)，所以 \(c=\sqrt{5^2-4^2}=3\). 焦点在长轴上且关于原点对称，故焦点坐标为 \(F_1\left(-3,0\right)\)、\(F_2\left(3,0\right)\).
\end{solution}

\begin{problem}
如右图所示的三视图是什么几何体？

\bitmapfigure[max width=6.0cm]{img/1977/shandong_a/q6_fig1.png}
\end{problem}

\begin{solution}
由俯视图可见，上、下底面都是正方形，且两个正方形的中心和对应边方向相同；由主视图、左视图可见，侧面轮廓都是上窄下宽的等腰梯形. 因此该三视图表示的几何体是正四棱台.
\end{solution}

\begin{problem}
设 \(AD\) 是 \(\triangle ABC\) 的高，在线段 \(BC\) 上取点 \(E\)，使 \(\angle EAC=\angle BAD\)，延长 \(AE\) 交 \(\triangle ABC\) 外接圆于点 \(F\). 求证：\(AF\) 是外接圆的直径.

\bitmapfigure[max width=4.9cm]{img/1977/shandong_a/q7_fig1.png}
\end{problem}

\begin{solution}
连接 \(FC\). 因为 \(A\)，\(E\)，\(F\) 三点共线，所以 \(\angle FAC=\angle EAC=\angle BAD\). 又因为 \(A\)，\(B\)，\(C\)，\(F\) 四点共圆，点 \(B\) 与点 \(F\) 在弦 \(AC\) 的同侧，故它们所对的同一弦 \(AC\) 的圆周角相等，即 \(\angle AFC=\angle ABC\).

由于 \(AD\) 是 \(\triangle ABC\) 的高，且 \(D\) 在 \(BC\) 上，所以 \(AD\perp BC\)，从而 \(\angle ADB=90^\circ\). 在直角三角形 \(ABD\) 中，\(\angle ABD=\angle ABC\)，故 \(\angle BAD+\angle ABC=90^\circ\). 因此在 \(\triangle ACF\) 中，
\[
\angle ACF=180^\circ-\angle FAC-\angle AFC=180^\circ-\angle BAD-\angle ABC=90^\circ
\]
于是 \(\angle ACF\) 是外接圆所对弦 \(AF\) 的圆周角，且为直角. 根据圆周角所对的弦为直径，\(AF\) 是该外接圆的直径.
\end{solution}

\begin{problem}
\(x^2-x\cos \alpha+\frac{1}{4}\sin^2 \alpha=0\)（\(0<\alpha<\frac{\pi}{2}\)）为关于 \(x\) 的二次方程.

\begin{enumerate}
\item 当 \(\alpha\) 为何值时，方程有两个不相等的实数根？
\item 当方程有两个不相等的实数根时，它的根是什么？
\end{enumerate}
\end{problem}

\begin{solution}
\begin{enumerate}
\item 该方程关于 \(x\) 的二次项系数为 \(1\)，一次项系数为 \(-\cos\alpha\)，常数项为 \(\frac14\sin^2\alpha\). 因此判别式为
\[
\Delta=(-\cos\alpha)^2-4\cdot1\cdot\frac14\sin^2\alpha=\cos^2\alpha-\sin^2\alpha=\cos 2\alpha
\]
方程有两个不相等的实数根当且仅当 \(\Delta>0\)，即 \(\cos 2\alpha>0\). 由 \(0<\alpha<\frac\pi2\) 得 \(0<2\alpha<\pi\)，在此范围内 \(\cos 2\alpha>0\) 等价于 \(0<2\alpha<\frac\pi2\)，所以 \(0<\alpha<\frac\pi4\).
\item 当 \(0<\alpha<\frac\pi4\) 时，\(\Delta=\cos 2\alpha>0\). 根据一元二次方程求根公式，
\[
x=\frac{\cos\alpha\pm\sqrt{\cos 2\alpha}}{2}
\]
即两个根为 \(x_1=\frac{\cos\alpha+\sqrt{\cos 2\alpha}}{2}\) 和 \(x_2=\frac{\cos\alpha-\sqrt{\cos 2\alpha}}{2}\).
\end{enumerate}
\end{solution}

\begin{problem}
在半径为 \(R\) 的一块圆形铁板上，剪去一个圆心角为 \(\alpha\) 弧度的扇形，用剩下的铁板做一个圆锥形容器（铁板厚度和加工量不计）.

\begin{enumerate}
\item 用 \(R\) 和 \(\alpha\) 表示圆锥底面半径 \(r\)、圆锥高 \(h\)、圆锥容积 \(V\)；
\item 当 \(R=20\mathrm{~cm}\)，\(\alpha=\frac{2\pi}{5}\) 弧度时，计算圆锥容积 \(V\)（取 \(\pi=3.14\)，精确到 \(1\mathrm{~cm^3}\)）.
\end{enumerate}
\end{problem}

\begin{solution}
\begin{enumerate}
\item 剪去扇形后，剩余扇形的弧长成为圆锥底面圆的周长，而剩余扇形的半径 \(R\) 成为圆锥的母线长，因此：
\[
2\pi r=R\left( 2\pi-\alpha \right)
\]
所以 \(r=\frac{2\pi-\alpha}{2\pi}R\).
圆锥的母线、底面半径和高构成直角三角形，所以
\[
h^2=R^2-r^2=R^2\left[1-\frac{\left( 2\pi-\alpha \right)^2}{4\pi^2}\right]=\frac{4\pi\alpha-\alpha^2}{4\pi^2}R^2
\]
由于高度取正值，
\[
h=\frac{\sqrt{4\pi\alpha-\alpha^2}}{2\pi}R
\]
于是
\[
V=\frac13\pi r^2h=\frac{R^3\left( 2\pi-\alpha \right)^2}{24\pi^2}\sqrt{4\pi\alpha-\alpha^2}
\]
\item 代入 \(R=20\)、\(\alpha=\frac{2\pi}{5}\)，先化简半径和高：
\(r=\frac{2\pi-\frac{2\pi}{5}}{2\pi}\cdot20=16\)，\(h=\sqrt{20^2-16^2}=12\).
因此
\[
V=\frac13\pi\cdot16^2\cdot12=1024\pi
\]
取 \(\pi=3.14\)，得 \(V=1024\times3.14=3215.36\mathrm{~cm^3}\)，精确到 \(1\mathrm{~cm^3}\) 为 \(3215\mathrm{~cm^3}\).
\end{enumerate}
\end{solution}

\begin{problem}
我炮兵阵地位于 \(A\)，两观察所分别设于 \(C\)、\(D\)，已知 \(\triangle ACD\) 为一正三角形，且 \(CD=a\)，当目标出现于 \(B\) 时，测得 \(\angle CDB=45^\circ\)、\(\angle BCD=75^\circ\)，求炮目距离 \(AB\).

\bitmapfigure[max width=6.0cm]{img/1977/shandong_a/q9_fig1.png}
\end{problem}

\begin{solution}
在 \(\triangle BCD\) 中，\(\angle CDB=45^\circ\)、\(\angle BCD=75^\circ\)，所以
\(\angle CBD=180^\circ-45^\circ-75^\circ=60^\circ\). 由正弦定理，\(CD\) 对应 \(60^\circ\)，\(BC\) 对应 \(45^\circ\)，故
\[
BC=\frac{a\sin45^\circ}{\sin60^\circ}=\frac{a\cdot\frac{\sqrt2}{2}}{\frac{\sqrt3}{2}}=\frac{\sqrt6}{3}a
\]
因为 \(\triangle ACD\) 为正三角形，所以 \(AC=CD=a\)，且 \(\angle ACD=60^\circ\). 由图中各射线的位置，\(\angle ACB=\angle ACD+\angle DCB=60^\circ+75^\circ=135^\circ\). 在 \(\triangle ABC\) 中使用余弦定理，
\[
AB^2=a^2+\left(\frac{\sqrt6}{3}a\right)^2-2a\cdot\frac{\sqrt6}{3}a\cos135^\circ=\frac{5+2\sqrt3}{3}a^2
\]
由于 \(AB>0\)，所以炮目距离为
\[
AB=\sqrt{\frac{5+2\sqrt3}{3}}a=\frac{\sqrt{15+6\sqrt3}}{3}a
\]
\end{solution}

\begin{problem}
现有可建猪圈 \(L\) 米的材料，要建彼此相连的 \(m\) 间面积相等的矩形猪圈（如下图）（门在内，墙的厚度不计）.

\begin{enumerate}
\item 怎样设计每间猪圈的长和宽，才能使所建的猪圈的总面积最大？
\item 按这种设计，\(m\) 间猪圈长的总和与宽的总和有什么关系？
\end{enumerate}

\bitmapfigure[max width=5.8cm]{img/1977/shandong_a/q10_fig1.png}
\end{problem}

\begin{solution}
\begin{enumerate}
\item 设图中每排横向围墙的长度为 \(x\) 米，设每条纵向围墙的长度为 \(y\) 米. 由于共有上下两排横向围墙和 \(m+1\) 条纵向围墙，材料长度满足
\(2x+(m+1)y=L\)，所以 \(y=\frac{L-2x}{m+1}\)，且必须有 \(0<x<\frac L2\).

每间猪圈的长为 \(\frac{x}{m}\)，宽为 \(y\)，故总面积为
\[
S=m\cdot\frac{x}{m}\cdot y=xy=\frac{x(L-2x)}{m+1}
\]
配方得
\[
S=\frac{L^2}{8(m+1)}-\frac{2}{m+1}\left( x-\frac L4 \right)^2\le \frac{L^2}{8(m+1)}
\]
等号成立当且仅当 \(x=\frac L4\). 此时
\[
y=\frac{L-2x}{m+1}=\frac{L}{2(m+1)}
\]
每间猪圈的长为 \(\frac{x}{m}=\frac{L}{4m}\).
所以每间猪圈的长为 \(\frac{L}{4m}\) 米、宽为 \(\frac{L}{2(m+1)}\) 米时，总面积最大，最大总面积为 \(\frac{L^2}{8(m+1)}\) 平方米.
\item 按图中所有围墙的长、宽统计，\(m\) 间猪圈的长边共有上下两排，长度总和为 \(2x=\frac L2\)；宽边共有 \(m+1\) 条，长度总和为 \((m+1)y=L-2x=\frac L2\). 因而长的总和与宽的总和相等.
\end{enumerate}
\end{solution}

\section{参考题}

\begin{problem}
求极限：\(\displaystyle\lim_{x\to\infty}\frac{\e^x}{x^2}\)
\end{problem}

\begin{solution}
当 \(x\to+\infty\) 时，分子 \(\e^x\to+\infty\)，分母 \(x^2\to+\infty\)，构成 \(\frac{\infty}{\infty}\) 型；在充分大的 \(x\) 上，分母及其导数均不为零，且分子、分母都可导，因此可使用洛必达法则. 第一次使用得
\[
\lim_{x\to\infty}\frac{\e^x}{x^2}=\lim_{x\to\infty}\frac{\e^x}{2x}
\]
右端仍为 \(\frac{\infty}{\infty}\) 型，再次使用洛必达法则得
\[
\lim_{x\to\infty}\frac{\e^x}{2x}=\lim_{x\to\infty}\frac{\e^x}{2}=+\infty
\]
所以
\(\displaystyle\lim_{x\to\infty}\frac{\e^x}{x^2}=+\infty\)，即该函数在 \(x\to+\infty\) 时无有限的实数极限.
\end{solution}

\begin{problem}
利用定积分证明：椭圆 \(\frac{x^2}{a^2}+\frac{y^2}{b^2}=1\) 的面积 \(S=\pi ab\).
\end{problem}

\begin{solution}
设 \(a>0\)、\(b>0\). 椭圆在第一象限的上半支满足
\(y=b\sqrt{1-\frac{x^2}{a^2}}\)，所以第一象限面积为
\[
S_1=\int_0^a b\sqrt{1-\frac{x^2}{a^2}}\,\mathrm{d}x=\frac ba\int_0^a\sqrt{a^2-x^2}\,\mathrm{d}x
\]
令 \(x=a\sin t\)，则 \(\mathrm{d}x=a\cos t\,\mathrm{d}t\)，且当 \(x=0\) 时 \(t=0\)，当 \(x=a\) 时 \(t=\frac\pi2\). 在 \(0\le t\le\frac\pi2\) 内 \(\cos t\ge0\)，故 \(\sqrt{a^2-x^2}=a\cos t\). 因此
\[
S_1=ab\int_0^{\frac\pi2}\cos^2t\,\mathrm{d}t=ab\int_0^{\frac\pi2}\frac{1+\cos2t}{2}\,\mathrm{d}t=ab\left(\frac t2+\frac{\sin2t}{4}\right)\bigg|_0^{\frac\pi2}=\frac{\pi ab}{4}
\]
椭圆关于 \(x\) 轴、\(y\) 轴均对称，四个象限面积相等，所以
\(S=4S_1=\pi ab\). 这就证明了椭圆面积公式.
\end{solution}
